\chapter*{译校说明}
\markboth{译校说明}{译校说明}
\addcontentsline{toc}{chapter}{译校说明}

\section*{一、底本与体例}
本书译自 Raf Bocklandt 所著 \textit{A Gentle Introduction to Homological Mirror Symmetry}（London Mathematical Society Student Texts 99，Cambridge University Press，2021）。译文按原书三部分、十二章的顺序排布；前言、正文、参考文献与索引齐备。参考文献保留英文原文，以便读者按原条目检索；正文中的方括号数字（如 [176]）为原书文献编号，可点击跳转至文献表。索引按中文术语的拼音排序，条目页码亦可点击跳转。

\section*{二、术语约定}
术语以中文学界与物理学界的通行译法为准，全书统一。需要特别说明的有：\textit{quiver} 译“箭图”，\textit{sheaf} 译“层”，\textit{scheme} 译“概形”，\textit{stack} 译“栈”，\textit{toric variety} 译“环簇”（\textit{toric divisor} 作“环簇除子”，\textit{toric geometry} 作“环簇几何”），\textit{ample line bundle} 与 \textit{very ample line bundle} 分别译“丰富线丛”“极丰富线丛”，\textit{twisted complex} 译“扭曲复形”，\textit{grading} 译“分次”，\textit{line field} 译“线场”，\textit{vanishing cycle} 译“消失圈”，\textit{exact}（辛几何语境）译“正合”（正合辛流形、正合拉格朗日子流形、正合 Fukaya 范畴）。人名、定理名（Morita、Koszul、Morse、Floer、Fukaya、Bridgeland、King、Landau--Ginzburg、Calabi--Yau 等）不译。

\subsection*{保留原文的专业术语}
同调镜像对称的文献长期以英文流通，若干术语在中文里尚无稳定、无歧义的通行译名；强行直译反而妨碍读者对照原始文献。因此本书对下列术语一律保留英文，首次出现处或本说明中给出中文解释：
\begin{itemize}
  \item \textit{wrapped}（\textit{wrapped Fukaya category}、\textit{wrapped Floer complex}、\textit{partially wrapped}）：按 Reeb 流“卷绕”无穷远端的构造，旧译“卷绕”易与曲线的卷绕数混淆，故保留英文；
  \item \textit{thimble}（\textit{Lefschetz thimble}）：Lefschetz 纤维化中由消失圈扫出的拉格朗日“顶针”，与 \textit{Lefschetz pencil}（莱夫希茨铅笔，另一概念）无关，为免混淆保留英文；
  \item \textit{stop}（\textit{partially wrapped Fukaya category} 的 stop）：截断 Reeb 流、把柱状端“堵住”的超曲面，旧译“停止子”易被误读为算子，故保留英文；
  \item \textit{Fukaya}（Fukaya 范畴）、\textit{Floer}、\textit{Reeb}、\textit{Lefschetz}、\textit{Landau--Ginzburg}、\textit{Novikov}、\textit{Kontsevich}、\textit{SYZ} 等专名：不译。
\end{itemize}

\section*{三、校订标记}
本书在忠实翻译之外，另有三类校订标记与一类导读文字，均与原书正文区分：
\begin{enumerate}
  \item \textbf{本节重点}：每节开头的浅色方框为译校者所加，提示该节的核心对象、所用工具与主要结论；原书无此内容。若只需原书内容，可将 \texttt{main.tex} 中的 \verb|\showtakeawaytrue| 改为 \verb|\showtakeawayfalse|，即可一次性隐藏全部方框，不影响其余排版。
  \item \textbf{章末评注}：每章末尾的“本章回顾与评注”方框为译校者所加，用于串讲本章知识脉络、复述核心结论、辨析易错之处，并就该章的处理方式与在全书中承上启下的作用作简要评论；原书无此内容。同样可在 \texttt{main.tex} 中用 \verb|\showchapterreviewfalse| 整体隐藏。
  \item \textbf{图题}：原书插图多为无题示意图，本书所加图题系译校者据图意拟写，仅作阅读辅助；每个图题后均标注“（原书 p.$n$）”，给出该图在原书中的页码，便于对照。
  \item \textbf{校订批注}：正文中的 LaTeX 注释（\verb|% [Tutor Note: ...]|）与书末“附录：校订批注”同文对应。批注分三类——易错预警（符号约定、指标升降、张量缩并、代数变换中容易出错处）、含糊点拨（原书以“显然”“容易看出”跳过论证处，说明省略了什么、依据哪条定理）、背景直觉（构造的几何或物理含义）。凡原书疑似笔误，均在批注中如实说明；正文按上下文处理并在批注中标出。
\end{enumerate}

\section*{四、责任声明}
译文力求准确，但同调镜像对称涉及辛几何、代数几何与表示论的交叉，术语与符号约定在不同文献间差异很大；凡是与原书不符之处，应以原书为准。翻译与校订中的错误由译校者负责，欢迎读者指正。
